% Ad Familiares 10.17 --- headnote
% ad-familiares-10-17, 19 May 43 BC, on the march toward Lepidus (Transalpine Gaul)

\textit{L. Munatius Plancus to Cicero, written on the march to
join Lepidus around the 18th or 19th of May 43 BC --- Perseus
dateline \textit{Scr. in itinere ad Lepidum xiv aut x iiii K.
Iun., a. 711 (43)}. The dispatch reports a tightening
situation in Transalpine Gaul after Mutina: Antony, with his
advance units, has reached Forum Iulii (Fr\'ejus) on the Ides
of May, with Ventidius two days behind him; Lepidus has
established his camp at Forum Voconii, twenty-four miles
inland, and has written to Plancus that he will wait there
for him. The whole campaign now turns on whether Lepidus is
in fact going to fight Antony, hold his army together, and
let Plancus's legions march in --- or whether he will be
absorbed.}

\textit{The second section is the one Plancus needed to write
in his own voice: an apology for sending his brother L.
Plotius Plancus, recovering from a near-fatal exhaustion,
back to Rome rather than keeping him at his side in camp.
With the consuls Hirtius and Pansa both dead from the
Mutina fighting, the city has been left ``stripped bare''
of senior magistrates, and the urban praetor's chair needs
a citizen of weight; Plancus has reasoned that his sick
brother will serve the state better at Rome than in the
field. He asks the Senate to lay any blame for the decision
on his own want of foresight, not on his brother's loyalty.
The third section reports that Lepidus has finally sent
him a hostage for good faith, Apella, and commends to
Cicero L. Gellius, his go-between in those negotiations
--- whose name carries a daggered crux (``de tribus
fratribus Segaviano''), an uncertain phrase from a corrupt
passage in the manuscripts.}

\textit{The dispatch is the immediate forerunner of 10.18 (the
report from the field of his own march toward Forum Voconii)
and shows the cautious confidence of a commander who still
believes Lepidus can be made to act: ``if he and fortune
together keep everything intact for me, I undertake to you
that I shall bring the business to a successful conclusion
swiftly.'' Within ten days Lepidus's army will have gone over
to Antony, and Plancus will be writing very different
letters.}
